\documentclass[12pt]{article}
\usepackage[utf8]{inputenc}
\usepackage[T1]{fontenc}
\usepackage{hyperref}
\title{Kominimung Language Resources for K-12 Curriculum Development}
\author{Compiled by Claude Sonnet 5.5 (claude-sonnet-5-5) for Andr\'as Kornai}
\date{}
\begin{document}
\maketitle

\section{Introduction}
Kominimung (ISO 639-3 xoi) is a little-documented language of Madang Province, Papua New Guinea. Reported speaker numbers are about 320 (2003) and 328 (1973 data, rated only 20 percent certain). Sources differ on its classification.

\section{Primary Educational Resources}
No educational materials were found.

\section{Linguistic and Reference Materials}
The search found no grammar or dictionary. It found only general reference pages.
\begin{itemize}
\item Wikipedia: \url{https://en.wikipedia.org/wiki/Kominimung_language}
\item Endangered Languages Project entry: \url{https://endangeredlanguages.com/lang/10697}
\item Endangered Languages Project bibliography: \url{https://endangeredlanguages.com/lang/10697/bibliography}
\end{itemize}

\section{Religious and Community Resources}
\begin{itemize}
\item Joshua Project people-group profile: \url{https://joshuaproject.net/people_groups/12763}
\end{itemize}

\section{Digital and Audio Resources}
None were found.

\section{Curriculum Development Context}
Kominimung has essentially no published teaching resources. Any curriculum would need new documentation work with the community. The bibliography page above is the best lead for existing survey material, such as Z'graggen's 1975 survey cited there.
\end{document}